\documentclass[a4paper,14pt]{extarticle}

\usepackage[left=30mm, right=10mm, top=20mm, bottom=20mm]{geometry}

\usepackage[T2A]{fontenc}
\usepackage[utf8]{inputenc}
\usepackage[russian]{babel}
\usepackage{tempora}
\usepackage{type1cm}
\usepackage{indentfirst}

\usepackage{minted}
\usepackage{xcolor}

\usepackage{siunitx}
\usepackage{booktabs}
\usepackage{multirow}


\usepackage{csquotes}
\usepackage[
    parentracker=true,
    backend=biber,
    style=gost-numeric,
    sorting=nty
]{biblatex}

\usepackage[nodisplayskipstretch]{setspace}
\usepackage[none]{hyphenat}
\usepackage{fancyhdr}
\usepackage{titlesec}
\usepackage{tocloft}

\usepackage{amssymb}
\usepackage{amsmath}
\usepackage{tikz}
\usepackage{pgfplots}
\usepackage{graphicx}
\usepackage{caption}
\usepackage{subcaption}
\usepackage[hidelinks]{hyperref}
\usepackage[russian]{cleveref}
\usepackage{bookmark}

\pgfplotsset{compat=1.18}

\onehalfspacing
\sloppy

\pagestyle{fancy}
\fancyhf{}
\fancyfoot[R]{\thepage}
\renewcommand{\headrulewidth}{0pt}

\definecolor{codebg}{rgb}{0.12,0.12,0.14}
\definecolor{codeframe}{rgb}{0.40,0.40,0.45}
\definecolor{keywordcolor}{rgb}{0.45,0.65,1.0}
\definecolor{stringcolor}{rgb}{1.0,0.55,0.45}
\definecolor{commentcolor}{rgb}{0.55,0.75,0.60}
\definecolor{textcolor}{rgb}{0.90,0.90,0.90}

\renewcommand{\theFancyVerbLine}{
	\sffamily\small\color{black}\arabic{FancyVerbLine}
}

\setminted{
    bgcolor=codebg,
	fillcolor=codebg,
	framesep=1em,
	frame=single,
    rulecolor=\color{codeframe},
    fontsize=\normalsize,
    baselinestretch=1.1,
    linenos,
    numbersep=6pt,
    tabsize=4,
    breaklines,
    breakbytoken,
	breaksymbolleft={},
    autogobble,
    showspaces=false,
    showtabs=false,
}

\usemintedstyle{lightbulb}
\setmintedinline{
    bgcolor=codebg
}

\makeatletter
\setlength{\@fptop}{0pt}
\setlength{\@fpsep}{8pt}
\setlength{\@fpbot}{0pt plus 1fil}
\makeatother

\newcommand{\circledequals}{\mathbin{\tikz \node[draw,circle,inner sep=1pt] {$=$};}}

\DeclareMathOperator{\sgn}{sgn}

\begin{document}
    \begin{titlepage}
		\centering
		
		\includegraphics{images/FEFU-logo.jpg}\\
		МИНИСТЕРСТВО ОБРАЗОВАНИЯ И НАУКИ РОССИЙСКОЙ ФЕДЕРАЦИИ\\
		Федеральное государственное автономное образовательное учреждение высшего образования\\
		\textbf{«Дальневосточный федеральный университет»}

		\rule{\linewidth}{0.8mm}\\[0.3cm]
		
		\textbf{ИНСТИТУТ МАТЕМАТИКИ И КОМПЬЮТЕРНЫХ ТЕХНОЛОГИЙ}\\
		\textbf{Кафедра математического и компьютерного моделирования}\\
		\vspace*{0.9cm}
		
		Кулахсзян Сергей Грайрович\\
        \vspace*{0.3cm}

		ЛАБОРАТОРНАЯ РАБОТА № 1 \\
        ЗАДАЧА № 8 \\
		\vspace*{0.5cm}

		\textbf{Задачи на площадь криволинейной трапеции}\\
		\vspace*{0.5cm}

		Направление подготовки 02.03.01сцт Сквозные цифровые технологии,\\ 
        бакалавриатская программа «Математика и компьютерные науки»\\ 
        Очной формы обучения\\
		\vspace{1.5cm}

		
		\vfill
		\centering
		г. Владивосток\\
		2026
	\end{titlepage}

    \tableofcontents
    \thispagestyle{empty}

	\newpage
    \pagenumbering{arabic}
    \setcounter{page}{3}
    \section{Постановка задачи}
    
\end{document}